\documentclass[nocoverpage,swedish,g5paper]{thesis}
%
%   optional options to documentclass:
%
%   coverpage   : Create both cover, inside front and text.
%                 Useful for web publishing.
% 
%   nocoverpage : Inner part of thesis only, do not create cover sheet.
%                 Useful for printing.
%   
%   onlycoverpage : Only create cover page. Ignores all text.
%                   Useful for printing.  
%
%   onlytext : Only print the text of the work. No cover and no inside front.
%              Useful for proof-reading copies.
%  
%  g5paper, s5paper, a4paper : Choose paper format, 
%
%  9pt, 10pt, 11pt, 12pt : Choose typeface size.
%
%  draft, final : Draft marks errors with a black box in text.
%
%  openright, openany : openright makes chapters only open at
%                       right hand pages.
%
%  * : Anything else is intepreted as the babel name of a
%      foregin language which is applied to the 'foregincommand'.
%
%
%  Default : s5paper,10pt,final,openright
%
%
%
%  required parameters
%
\title{Monte Carlo study of criticality in finite lattice models}
\author{Simon Sandell}
\date{June 2018}
\shortdate{2018}
\type{Master Thesis}
\department{Department of Theoretical Physics,\\School of Engineering Sciences}
\address{SE-106 91 Stockholm, Sweden}
\city{Stockholm}
\country{Sweden}
\publisher{Printed in Sweden by Universitetsservice US AB, Stockholm January 2010}
\copyrightline{\copyright\ Simon Sandell, June 2018}
\trita{FYS-9999:99}
\isbn{999-99-9999-999-9}
\issn{9999-999X}
\isrn{KTH/FYS/-{}-99:99-{}-SE}
\comment{Scientific thesis for the degree of Master in Science of Engineering (MSc) in the subject area of Theoretical physics.\\ \\ \textbf{Cover illustration:} Non-existent.}
%
%  optional parameters
%
\cplogo{\includegraphics[height=2.5cm]{./kthlogo.eps}}
\innerlogo{\includegraphics[height=2.5cm]{./kthlogo.eps}}
\subtitle{An investigation of the discrepancy in measured and calculated critical exponents of the lambda-transition in $^{4}He$, as well as calculation of the critical exponents of the transition in the 3 dimensional Ising-model (because why not?).}
\division{Condensed Matter Theory}
\centercomment{\centerline{Typeset in \LaTeX}}
%\foregincomment{Akademisk avhandling f\"or avl\"aggande av teknologie licentiatexamen (TeknL) inom
%\"amnesomr{\aa}det teoretisk fysik.}
%\dedication{To Someone}

%\usepackage[swedish]{babel}
\usepackage{verbatim}
\usepackage{bm}
\usepackage{enumerate}
\usepackage[dvips]{graphicx}
\usepackage{amsmath,amsfonts,amssymb}
\usepackage[utf8]{inputenc}
\usepackage{url}
\usepackage[square, comma, sort&compress,numbers]{natbib}

\unitlength=1mm

\def\slc#1{\setbox0=\hbox{$#1$}           % set a box for #1
    \dimen0=\wd0                                 % and get its size
    \setbox1=\hbox{/} \dimen1=\wd1               % get size of /
    \ifdim\dimen0>\dimen1                        % #1 is bigger
       \rlap{\hbox to \dimen0{\hfil/\hfil}}      % so center / in box
       #1                                        % and print #1
    \else                                        % / is bigger
       \rlap{\hbox to \dimen1{\hfil$#1$\hfil}}   % so center #1
       /                                         % and print /
    \fi}

\newcommand{\todo}[1]{(\textbf{TODO:} #1)}
\newcommand{\ud}{\mathrm{d}}
\newcommand{\dd}[2]{\frac{{\rm d}#1}{{\rm d}#2}}
\newcommand{\citeb}[1]{[\citen{#1}]}
\newcommand{\Ref}{[{\bf REF}]}
\newcommand{\Fig}{[{\bf FIG}]}
\newcommand{\chk}{[{\bf CHECK}]}
\newcommand{\im}{\mathrm{i}}
\newcommand{\Mpl}{M_{\rm Pl}}
\newcommand{\Mpr}{\bar{M}_{\rm Pl}}
\newcommand{\Ms}{M_*}
\newcommand{\Msr}{\bar{M}_*}
\newcommand{\ie}{{\it i.e.}}
\newcommand{\eg}{{\it e.g.}}
\newcommand{\hc}{{\rm h.c.}}
\newcommand{\trm}[1]{\textrm{#1}}
\newcommand{\oxx}{\omega_{XX'}}
\newcommand{\hamxy}{\sum\limits_{\langle i,j\rangle}\bm{s_i}\cdot\bm{s_j}}

\begin{document}

%\def\@cite#1{[#1]}

\begin{abstract}
Put abstract here.

And here.

And here..
\\\noindent \strut \\
{\bf Key words}: Monte Carlo, Phase Transition, Finite Size Scaling, 3D Ising model, 3DXY model,Critical Exponents,Universality
\end{abstract}

%\begin{otherlanguage}{swedish}
%\begin{foreginabstract}
%\todo{Skriv ett abstract p{\aa} svenska.}
%\\\noindent \strut \\
%{\bf Nyckelord}: Extradimensionella kvantf{\"a}ltteorier, universella extra dimensioner, ADD-modeller, Kaluza--Klein-m\"ork materia, neutrinomassor, LHC-fenomenologi
%\end{foreginabstract}
%\end{otherlanguage}

\begin{preface}
% !TEX root = thesis.tex
%        File: preface.tex
%     Created: mån feb 12 01:00  2018 C
% Last Change: mån feb 12 01:00  2018 C
%
\section{Numerical constants}
Throughout this thesis we let the following phyiscal quantities have the following values:
\begin{equation}
  k_B = 1
  \label{}
\end{equation}

\section{Shorthand notations and definitions}
\begin{table}[htpb]
  \begin{center}
    \begin{tabular}{l l}
      Notation & Definition\\
      MC & Monte Carlo\\
      FSS & Finite Size Scaling\\
      $L$ & Linear system size measured in units of lattice spacing\\
      $\beta \equiv \frac{1}{k_B T}$ & Inverse temperature\\
    \end{tabular}
  \end{center}
  \caption{Clarification of notation}
\end{table}



\end{preface}

\tableofcontents

% This separates the introduction from the main part of the thesis.
\mainmatter


\chapter{Introduction}
%        File: preface.tex
%     Created: mån feb 12 01:00  2018 C
% Last Change: mån feb 12 01:00  2018 C
%
% TODO: Write something interesting about ising3d and new scaling method.

Phase transitions are interesting phenomenae. They are part of everyday life, with such manifestations as ice melting, water boiling, etc, but there are endlessly many more examples in nature, for example superfluid and superconducting phaste transitions. In any system with bulk properties experiencing a discontinous or  undergoing a change of order. The best and most accurate theory describing phase transitions is called Renormalization Group-theory (RG) and has been very successfull in explaining critical phenomenae, but the precise behaviour of physical quantities near and at the transtition temperature where the theory makes it's claims is for many systems hard to both calculate from the theory and measure in experiments.

The superfluid transition in helium got it's nickname, ``The lambda transition'', due to the peculiar shape of the graph of heat capacity versus temperature, which resemples the greek letter. The transition happens at about 2.2 K, at normal pressure, where the liquid transitions into a superfluid state. It is important as a experimental verification of the Renormalization Group-theory (RG), as it is one of few systems where experimental measurement of these quantities is feasible. 

The lambda phase transition in Helium-4 has been studied extensively both in theory, simulations and experiments. In 2003 the results of measurements performed in orbit was published\cite{Lipa2003} stating a value for the critical exponent alpha to unprecedented precision.

The experiment was performed on the Space Shuttle Columbia, the reason for this is that a gravitational gradient will introduce pressure, causing a transition broadening. In theory, the heat capacity should diverge to infinity at the transition temperature.
The measurements of the heat capacity as well as the superfluid density was performed, to temperatures in the range of $nK$ of $T_c$, with no evidence of broadening.

Theoretical calculations of Sokolov and Nikitina \cite{Sokolov2016} have given predictions of the critical exponents which to the stated precision agrees with the value seen in experiment.
However, numerical studies performed by Campostrini et. al. \cite{Campostrini2006} has predicted values that are incompatible with both experiment and the theoretical predictions. 

The aim of this thesis was to achieve numerical estimations to a higher precision than previous studies by virtue of a new approach to finite size scaling corrections. The method should work for numerous systems, we chose to investigate the 3DXY model since it is of interest to see if discrepancies will dissapear with higher precision, and also the 3D Ising-model, since it is also a very important model, and it's implementation is very similar to that of the 3DXY-model. 

Our results show a better/the same/worse agreement with the experimental value. Agreement shows that numerical studies are a valid method of study of these systems, and discredits the idea that numberical simulations may be part of a different but closely related universality class.




\chapter{Background}\label{ch:Background}
%        File: preface.tex
%     Created: mån feb 12 01:00  2018 C
% Last Change: mån feb 12 01:00  2018 C
%
% TODO: Scaling and universality: more equations.
% scaling laws and definitions of critical exponents, maybe remove? maybe write some text. move to appendix or other place?
% There should be some parts about xy model and ising3d model here.
In this section we will describe the models we simulate using Monte Carlo, as well as establish some theoretical background to explain how physical information can be extracted from simulated systems with sizes $10^X$ times smaller than actual realizations of the transition studied in experiments.
\section{Models} % define equations, define quantities.
The goal of this thesis is to estimate the critical exponents of the lambda phase transition (\lpt) in $^4$He by Monte Carlo (MC) simulation of a model. 
The principle of Universality allows us to freely chose a model for our simulation as long as it is part of the same Universality Class as the \lpt, since all those models are guaranteed to have the exact same critical properties. 
The three criteria that define universality classes are the symmetry group of the order parameter, the physical dimensions and the range of the interactions.
The class of the lambda transition in $^4$He is represented by the symmetry group $O(2)$, short range interactions and $d=3$. With $d=2$ instead the class changes to that of the famous topological Kosterlitz-Thouless transition.

The model we use in our simulation is the so-called XY-model in 3 dimension and allowing only nearest neighbour interactions on a cubic lattice.
We employ periodic boundary conditions, so that the nearest neighbours of any given lattice site $\bm{s} = (x,y,z)$ is defined as 
\begin{align}
    \bm{n_{1,2}} &= (\trm{mod}(L+x\pm1), &   &y,                    &    z)&\\
    \bm{n_{3,4}} &= (x,                  &   &\trm{mod}(L+y\pm1),   &    z)&\\
    \bm{n_{5,6}} &= (x,                  &   &y,                    &    \trm{mod}(L+z\pm1))&
\end{align}
where the lattice indices $(x,y,z)$ range from $[0,L-1]$.

The Hamiltonian is written as 
\begin{equation}
  \label{xyham}
  H_{\trm{3DXY}} = -K\sum\limits_{\langle i,j\rangle} \bm{s_i}\cdot\bm{s_j} = -K\sum\limits_{\langle i,j\rangle} \cos(\alpha_i - \alpha_j)
\end{equation}
where $K$ is the coupling constant, which we during the rest of the thesis let $K=1$.
In the XY-model, the spins are confined to a plane, so they can be parametrized by $ \bm{s_1} = (\cos(\alpha_i),\sin(\alpha_i)$, hence the second equality in \ref{xyham}.

Next, we define some of the quantities which we calculate during simulation.
We define the magnetization as 
\begin{equation}
  M = \left| \sum_i \bm{s_i}\right|
  \label{}
\end{equation},
the energy as 
\begin{equation}
  E = -\sum\limits_{\langle i,j \rangle} \cos(\alpha_i - \alpha_j)
  \label{}
\end{equation},
the superfluid stiffness/helicity modulus as 
\begin{equation}
  L\rho_s = -\langle E\rangle - \frac{1}{T}\left\langle\left(\sum_i \sin(\alpha_i - \alpha_{i+(1,0,0)})\right)^2\right\rangle
  \label{}
\end{equation}
The superfluid stiffness/density/helicity modulus is a quantity representing the shift in free energy introduced by a twist of the spins.
To improve statistical results, we also compute
\begin{align}
  L\rho_s = -\langle E\rangle - \frac{1}{T}\left\langle\left(\sum_i \sin(\alpha_i - \alpha_{i+(0,1,0)})\right)^2\right\rangle\\
  L\rho_s = -\langle E\rangle - \frac{1}{T}\left\langle\left(\sum_i \sin(\alpha_i - \alpha_{i+(0,0,1)})\right)^2\right\rangle
  \label{}
\end{align}
In the regular liquid phase, the energy of the system is unaffected by uniform boundary motion. The superfluid stiffness is a quantity which measures the shift in energy due to a shift of the angles along a boundary.
It is interesting as it scales as 

The helicity modulus $\Upsilon(T)$ is a measure of the system to a ``phase twisting'' field, for a superfluid system the helicity modulus is related to the superfluid density
$\rho_s = (m/\hbar)^2 \Upsilon(T)$.
Fischer et al showed in 1973 that this quantity can be calculated within the framework of equilibrium statistical mechanics, but one need to go beyond bulk properties.
Operational definition for spherical model (1952 Barber Fischer) and ideal Bose gas (Barber 1977) under periodic and antiperiodic boundary conditions:
\begin{equation}
  \Upsilon(T) = \trm{lim}_{L\rightarrow \infty} (2L^2/\pi^2) (F^{1/2} (T;L)-F^0(T;L)).
  \label{}
\end{equation}
Rudnick Jasnow (1977) operational definition: 
\begin{equation}
  \Upsilon(T) = \left. \frac{  \pa^2F(T;k_0)}{\pa k^2_0}\right|_{k_0 =0}
  \label{}
\end{equation}
Thought: the boundary conditions induce a long wavelength twist of the wavenumber $k_0$ of the order parameter.
Showed that Josephson relation ($\nu = 2\beta -\eta\nu = 2 - \alpha -2\nu = (d -2)\nu$) (1966) is exact using this definition and epsilon-expansion.
The superfluid density is in the thermodynamic limit a constant at $T_c$.



\section{Finite size scaling}
Consider a system with a set of coupling constants $[K]$ and linear finite size $L$. 
The singular part of the free energy scales as
\begin{equation}
  f_s([K],L^{-1}) = l^{-d}f_s([K],lL^{-1}).
  \label{}
\end{equation}
Close to a fixed point of the RG, we can write this equation in terms of right eigenvectors of the linearized RG-transform,
\begin{equation}
  f_s(t,h,K_3,\cdots,L^{-1}) = l^{-d}f_s(tl^{y_t},hl^{y_h},K_3 l^{y_3},\cdots,lL^{-1}).
  \label{}
\end{equation}
It is evident that the inverse size of the system is in fact a relevant eigenvector with eigenvalue $y_L = 1$, and only becomes irrelevant in the thermodynamic limit $L^{-1}\rightarrow 0$.
The models we study in this paper have the external field set to zero, if we let the scaling parameter $l = L$, we can write the scaling form of the free energy as
\begin{equation}
  f_s(t,L^{-1}) = L^{-d} F_f(t L^{1/\nu})
  \label{}
\end{equation}
where we used $\nu = 1/y_t$ and the scaling law $2 -\alpha = d\nu$.
$F_f(tL^{1/\nu})$ is the form function for the free energy.
From this equation we can derive the finite size scaling behaviour of any physical quantities by utilizing their relation to the free energy. 
For example, we can find the scaling form of the specific heat as
\begin{equation}
  c_V = \frac{\partial ^2 f_s}{\partial t^2} \sim L^{-d +2/\nu}F_{c}(tL^{1/\nu}) = L^{\alpha/\nu} F_{c}(tL^{1/\nu}).
  \label{}
\end{equation}
This form function will have a maximum at some shifted value of $T$, 
\begin{equation}
  T_c(L) = T_c +  a_0 L^{-1/\nu}.
  \label{}
\end{equation}


Similarly for the susceptibility,
\begin{align}
  \chi = \frac{\partial^2 f}{\partial h^2}
  \sim L^{-d +2y_h}F_{\chi}(tL^{1/\nu}) = L^{2-\eta}F_{\chi}(tL^{1/\nu})
  \label{}
\end{align}
where we used that $ 2(d-y_h) = d -2 +\eta$.
The magnetization goes as
\begin{equation}
  M = \frac{\partial f}{\partial h} \sim L^{-d + y_h} = L^{-\beta/\nu}
  \label{}
\end{equation}


\section{Scaling laws}
From renormalization-group theory we learn that there only exist two independent critical exponents from which all the others are derived.
This was not a surprise, in fact, it was foreshadowed by the many scaling laws proposed both from prior theory and experiments.
The framework of renormalization-group define the two independent exponents as the eigenvalues of rescaling operators.
Here we state the various relations between exponents and their definitions in terms of limiting behaviour of physical quantities.
\begin{align}
  y_t &= \frac{1}{\nu}\\
  y_h &= \frac{2-\eta +d}{2}\\
  \Delta &= y_h/y_t\\
	2 - \eta &= \frac{\gamma}{\nu} = d\frac{\delta -1}{\delta+1}\\
	\nu d &= 2-\alpha = 2\beta + \gamma = \beta(\delta+1) = \gamma\frac{\delta +1}{\delta -1}
\end{align}
\section{Definitions of critical exponents}
For $h=0$,
\begin{align}
  f &\sim l^{-d}f(tl^{y_t},hl^{y_h},K_3l^{y_3},\cdots,L^{-1}l)\\
  C &\sim t^{-\alpha}\\
  M &\sim (-t)^{\beta}\\
  \chi &\sim |t|^{-\gamma}\\
  \xi &\sim |t|^{-\nu}\\
\end{align}  
For $t =0$,
\begin{align}
  h &\sim m^{\delta}\\
  \langle m(r)m(0)\rangle &\sim r^{-d +2 -\eta}
  \label{}
\end{align}


\chapter{Method}\label{ch:Method}
%        File: preface.tex
%     Created: mån feb 12 01:00  2018 C
% Last Change: mån feb 12 01:00  2018 C
%

\section{Monte Carlo}
In statistical physics, predictions come from thermal averages calculated from 
\begin{equation}
  \langle A\rangle = \frac{1}{Z}\trm{Tr} e^{-H/T} = \sum_x A(x) P(x)
\end{equation}
where  $P(x) = (1/Z)\trm{exp}(-H(x))$ is the Boltzmann distribution, and the sum is a sum over all possible states of the system.
Models of physical systems are described by writing down the mathematical expression for the Hamiltonian. However, an expression for the partition function $Z$ can in most cases not be found. So using the first representation to calculate averages is for most models not possible. The other form shows another way to calculate the averages, at least approximatesly, by simulation. 
Monte Carlo averages have the form 
\begin{equation}
  \langle A \rangle = \frac{1}{N} \sum_x A(x) \pm \frac{\sigma	}{\sqrt N}
\end{equation}
where the states $x$ are Boltzmann distributed.

%The Boltzmann distribution may be unsuitable for some systems/configurations we want to study. In practice, we can use any distribution we want, if we introduce the proper normalization.
%\begin{equation}
 % \langle A \rangle = \frac{\frac{1}{N}\sum_y \frac{A(y) e^{-H(x)/T}}{P'(y)}}{\frac{1}{N}\sum_y \frac{e^{-H(y)/T}}{P'(y)}}  
%\end{equation}
 %where the states $y$ are distributed according to $P'(y)$.

%Importance sampling is a way of reducing the error in the Monte Carlo estimate. Instead of choosing sample points distributed uniformly, one can try to choose the most probable states. Since thermal systems are Boltzmann-distributed, we want a way to randomly generate states distributed according to the Boltzmann distribution.

%//\section{Markov process}
%//A Markovian process is one for which the probility of getting to any state in the system is determined solely by the current state. I.e. the probability of going from state $x_i$ to $x_{j}$ is given by a transition probability $w(x_i \rightarrow x_j)$. The master equation describing the evolution of the probability distribution can then be written $P(x_{i+1}) = k
%// 
%//

\section{Histrogram extrapolation}
Histogram extrapolation is a method used to extrapolate data obtained from a simualtion.
By calculating additional averages during the simulation proceeds, one can obtain data for a range of temperatures surrounding the actual temperature the simulation is running at.
The averages are calculated as
\begin{equation}
  \langle A \rangle_{\beta_0} =  \frac{\sum_{x}A_x e^{-\beta_0 H_x}}{\sum_{x}e^{-\beta_0 H_x}}
\end{equation}
To get the average of A at some different temperature, say $\beta_1$, we can do
\begin{align}
  \langle A\rangle_{\beta_1} =  \frac{\sum_{x}A_x e^{-\beta_1 H_x}}{\sum_{x}e^{-\beta_1 H_x}} = \\
  = \frac{\sum_x A_x e^{-(\beta_1 - \beta_0)H_x} e^{-\beta_0 H_x}}{\sum_x e^{-(\beta_1 - \beta_0)H_x}e^{-\beta_0 H_x}} = \\
  = \frac{\langle A e^{-(\beta_1 - \beta_0)H}\rangle_{\beta_0}}{\langle e^{-(\beta_1 - \beta_0)H}\rangle_{\beta_0}}.
\end{align}
The range of temperatures to which extrapolation is viable will decrease with system size,
as loss in precision will occur when the exponential factor becomes large. To minimize the exponential factor, one can shift the energy of the system by a constant, which we call $H_{max}$. A constant shift in energy does not affect the values, 
\begin{equation}
  \langle A \rangle = \frac{e^{\beta H_{max}}}{e^{\beta H_{max}}}\frac{\sum_x A_x e^{-\beta H_x}}{\sum_x e^{-\beta H_x}}= \frac{\sum_x A_x e^{-\beta(H_x-H_{max})}}{\sum_x e^{-\beta(H_x-H_{max})}}.
  \label{}
\end{equation}
This technique can extend the range of temperate, but fluctuations in energy grow with the system size, so for high enough system sizes, extrapolation becomes fruitless.
The data gathered at the extrapolated temperature will be worse and worse the further we go, since we generate states that are probable for actual running temperature, but these states will likely become more and more unimportant for the extrapolated temperature.
\section{Wolff Algorithm}
The Wolff algorithm is a non-local update method, and is an improvement upon the Swendsen-Wang method in which spin bonds are scanned and either deleted or frozen. This will divide the lattice into clusters, which are then flipped/updated with certain probabilites. In the Wolff algorithm only a single cluster is generated, with the starting spin uniformly random. That way of initializing the cluster, one can expect to hit large clusters with higher probability, since probability of hitting a cluster should depend on its size.
Thus one can avoid the many single spin clusters usually generated in the Swendsen-Wang method.
For a markovian update method to give a chain of system configurations with distribution according to the Boltzmann-factor, it must satisfy two conditions, 

(i) for any given state of the simulated system, any other state must be reachable in a finite number of steps.

(ii) Non-periodicity, it should not be possible to return to a preivous state immediatley, only after a number of steps, $t =nk,~ n = 1,2,3\dots$.
This hold for the Wolff method, since a valid cluster consists of just shifting a single spin with a random angle, thus any configuration can be reached in $N \leq N_{\trm{spins}}$ steps. Also, since the first spin selected is always flipped, the condition of Non-periodicity is also fulfilled.
The algorithm must also have a transition probability configured so that after a long time, the desired stationary distribution is reached, in our case the Boltzmann-distribution.
Call the desired staionary distribution $\rho(X)$, the transition probability $T(X \rightarrow X')$ and the probability of state $X$ at markov timestep t $\rho(X,t)$.
Then the master equaion is 
\begin{equation}
  \rho(X, t+1) - \rho(X,t) = -\sum_{X'} T(X\rightarrow X')\rho(X,t) +\sum_{X} T(X'\rightarrow X)\rho(X',t).
\end{equation}
Then solving for the stationary solution, we get
\begin{equation}
  \sum_{X'}T(X\rightarrow X') \rho(X) =  \sum_{X}T(X'\rightarrow X) \rho(X')
\end{equation}
which is hard to solve in general, but with a very famous paritcular solution called the detailed balace solution can be found,
\begin{align}
  \frac{T(X\rightarrow X')}{T(X'\rightarrow X)} = \frac{\rho(X')}{\rho(X)} 
  \label{eq:detbal}
\end{align}
for all states $X,~X'$.

This solution can be decomposed as
\begin{equation}
  T(X\rightarrow X') = \omega_{XX'} A_{XX'}
\end{equation}
where $\oxx$ represents a trial probability and $A_{XX'}$ represents an acceptance probability.
We let $\omega_{XX'}$ satisfy
\begin{align}
  \oxx = \omega_{X'X}\\
  0 \leq \oxx \leq 1\\
  \sum_{X'} \oxx = 1.
\end{align}
Inserting this form of $T$ into the detailed balance equation then gives
\begin{equation}
  \frac{A_{XX'}}{A_{X'X}} = \frac{\rho_{X'}}{\rho_{X}}.
\end{equation}
In the metropolis algorithm one chooses the acceptance probability as follows
Step 1: From state $X$ propose a new trial state $X'$ with probability given by $\oxx$. Then accept that state with probability 
\begin{equation}
  A_{XX'} = 
  \begin{cases}
    	1 &\trm{ if } \frac{\rho(X')}{\rho(X)} \geq 1\\
	\frac{\rho(X')}{\rho(X)}&\trm{ if } \frac{\rho(X')}{\rho(X)} < 1 .
  \end{cases}
\end{equation}
The trial probability is usually just a uniform distribution.
\section{Wolff on the 3DXY-model}
In the 3DXY-model, each spin is characterized by a single value, their angle $\alpha$.
We define the hamiltonian of the 3DXY-model as 
\begin{equation}
  H = \beta\sum_{<s_i,s_j>} \cos(\alpha_i - \alpha_j)
\end{equation}
where $<s_i,s_>$ denote that only sites $s_i$ and $s_j$ which are nearest neighours should be included.
The probability of adding spin to the cluster is
\begin{equation}
 P(\alpha_i,\alpha_j,\alpha_u) = 1 - \exp(2\beta \cos(\alpha_j - \alpha_u)\cos(\alpha_i - \alpha_u))
  \label{pflip}
\end{equation}
where $\bm{u}$ is the normal to the plane through which the spins are reflected when added.
When a spin is added to the cluster, it is reflected through a plane.
The Wolff-algorithm is defined as follows for the 3DXY-model:
\begin{enumerate}[(i)]
\item Select a starting spin with uniform probability. Select a random plane with normal $\bm{u}$, with the angle $\alpha_u$  to reflect spins through. Reflect the starting spin through the plane, taking it to 
$$\alpha \rightarrow R(\alpha,\alpha_u) = \pi + 2\alpha_u - \alpha$$

and mark it as part of cluster.

\item Add it's neighbours to a list of perimeter spins, to be tried for inclusion in the cluster.

\item Pick out any element in the list. Try to add it to the cluster with probability as defined above. If it is accepted, add it to the cluster. Add nearest neighbours that are not part of the cluster already to the perimeter list.

\item Repeat (iii) until the perimeter list is empty.
\end{enumerate}
For the Wolff-algorithm, to show that it satisfies detailed balance, condiser a 
state $X$ and a state $X'$ differing by a flip of the cluster $c$. 
The transition probabilities obey
\begin{align}
  \frac{T(X\rightarrow X')}{T(X'\rightarrow X)} &= \prod_{\langle s_i s_j \rangle \in \bm{\partial} c} \frac{ 1- P(R(\alpha_i,\alpha_u),\alpha_j,\alpha_u)}{1- P(R(\alpha'_i,\alpha_u),\alpha'_j,\alpha_u)}
\end{align}
where the product is over nearest neighbour bonds where $s_i \in x, s_j \notin c$.
Thus we have that $ R(\alpha_i,\alpha_u) = \alpha'_i $, $R(\alpha'_i) = \alpha_i$ and $\alpha_j = \alpha'_j$.
We get
\begin{align}
  \frac{T(X\rightarrow X')}{T(X'\rightarrow X)}&= \prod_{\langle s_i s_j \rangle \in \bm{\partial} c} \frac{ \exp\left\{2\beta\cos(R(\alpha_i,\alpha_u) - \alpha_u)\cos(\alpha_j - \alpha_u)\right\}}{ \exp\left\{2\beta\cos(R(\alpha'_i,\alpha_u) - \alpha_u)\cos(\alpha'_j -\alpha_u)\right\}} = \\
  &= \prod_{\langle s_i s_j \rangle \in \bm{\partial} c} \frac{ \exp\left\{2\beta\cos(\alpha'_i - \alpha_u)\cos(\alpha'_j - \alpha_u)\right\}}{ \exp\left\{2\beta\cos(\alpha_i,\alpha_u) - \alpha_u)\cos(\alpha_j -\alpha_u)\right\}} = \\
  &= \frac{\rho(X')}{\rho(X)},
\end{align}
and so, detailed balance is satisfied.

\section{Model}
The XY-model and the Ising-model are both parth of the class of symmetry groups called $O(n)$, with $n=2,1$ respectively.
There are many different models in the universality class of the normal-superfluid Helimu-4 transition, the O(2)-symmetry group in three dimensions. Some examples are the $\phi^4$-model, the ddXY model etc.
In this thesis we have chosen the simulate the classical 3DXY-model. This is convenient because of the relative simplicity of the model, and that it is well known and studied.

The numerical method we used to perform the Monte-Carlo simulations was a Wolff algorithm written in C++. Since both the 3DXY-model and the 3D Ising model suffer from critcal slowing down, i.e. the equilibration time tends to infinity at the critical temperature, 
Both the 3DXY-model and the 3D Ising-model suffer from a phenomenon called critical slowing down. Near the transition temperature, the correlation length diverges, so thermal fluctuations propagate very long distances, and so the system takes a very long time to equilibrate to the equilibrum value.
One can wastly improve the quality-data versus cpu-time ratio by utilizing a global update algorithm such as the Wolff algorithm. 
\begin{equation}
  H_{\trm{3DXY}} = -K\sum\limits_{\langle i,j\rangle} \bm{s_i}\cdot\bm{s_j} = -K\sum\limits_{\langle i,j\rangle} \cos(\alpha_i - \alpha_j)
\end{equation}
We set $K=1$.
\begin{equation}
  Z = \trm{Tr} \exp\left[-\beta \left( H_\Omega - \int d^d \bm{r} H(\bm{r})\eta(\bm{r})\right)\right]
\end{equation}
\begin{equation}
  \label{}
\end{equation}


\chapter{Results}\label{ch:Results}
%        File: preface.tex
%     Created: mån feb 12 01:00  2018 C
% Last Change: mån feb 12 01:00  2018 C
%
% TODO: write some contaning text, include only relevant plots. explain results.
In this chapter we will first present the results of the equilibration studies. Then the main computational results will be presented, first the determination of the scaling correction $\omega$, the determinaion of the critical temperature and finally the estimations of the critical exponents $\eta$ and $\nu$.
\section{Equilibration}

\section{Previous results}
\begin{table}[!htpb]
\begin{center}
\begin{tabular}{l l l l l}
  Quantity 	& MC\cite{Campostrini2006} 	& pseudo-$\varepsilon$\cite{Sokolov2016} & Space\cite{Lipa2003}	& MC\cite{Jeon2012} \\
  $T_c$		& 2.20184(6)			& N/A					&			& 2.20186(9)\\
  $\alpha $	& -0.0151(3)			& -0.0117(31)				&-0.0127(3)		& -0.0124(30)\\
  $\nu	  $     & 0.6717(1) 			& 0.6706(10)				&0.6709(2)		& 0.6708(5)\\
  $\eta   $     & 0.0381(2) 			& 0.0376(28) 				&			& \\
  $\gamma $	& 1.3178(2) 			& 1.3159(8)				&			& \\
  $\beta  $     & 0.3486(1) 			& 0.3479(16)				&			& \\
  $\delta $	& 4.780(1) 			& 					&			& \\
  $\omega $     & 0.785(20) 			& N/A					&			& \\
  $\omega_2$ 	& 1.8(2) 			& N/A					&			& \\
\end{tabular}
\end{center}
\caption{3DXY-model }
\end{table}
Sokolov did pesudo-$\varepsilon$ expasion on 3dxy.\cite{Sokolov2016}
Lipa et al. did space stuff \cite{Lipa2003}
\begin{table}[!htpb]
\begin{center}
\begin{tabular}{l l l l}
  Quantity  & Hasenbusch & Conformal bootstrap& Our\\
  $T_c$	    & 4.511523(3)& N/A		& N/A\\
  $\alpha $ & N/A        & N/A		& N/A\\
  $\nu	  $ & 0.63002(10)& 0.629(1)	& N/A\\
  $\eta   $ & 0.03627(10)& 0.0341(5)	& N/A\\
  $\gamma $ & N/A        & N/A		& N/A\\
  $\delta $ & N/A        & N/A		& N/A\\
  $\beta  $ & N/A        & N/A		& N/A\\
  $\omega $ & N/A        & 0.80(1)	& N/A\\
  $\omega_2$& N/A        & N/A		& N/A\\
\end{tabular}
\end{center}
\caption{3D Ising-model }
\end{table}
Gliozzo used conformal bootstrap method on ising 3d.\cite{Gliozzi2014}
\subsection{Data Analysis}
The simulation program outputs raw data ordered in lines. Data from one ``simulation'' is contained in $N_{temperatures} \cdot N_{averages per simulation} \cdot N_{simulations}$.
To calculate quantities of interest from this raw data, a python script has been written, utilizing the numpy library. All available raw data is loaded into a numpy ndarray, and sorted by system size and temperature. Then for each block of data with one temperature and one system size, combined averages are calculated and from these, further functions of the averages are calculated, such as the Binder Parameter, it's derivative, Susceptibility, Superfluid density, etc. To estimate the statistical error in these estimations, we utilize a resampling method usually refered to as the jackknife method. For each block used in calculating the estimations, a jackknife function is called, which calculates the quantities again, but with a subblock of data omitted.
\subsection{3DXY-model}
Simulations were performed using the Wolff-cluster algorithm at the temperature $T_{run} = 2.2020~K$. Histogram extrapolation was performed to extrapolate to the range of temperatures $T_{range} = 2.2015-2.2030 ~ K$.
The simulations are structured so that a lattice is initialized in the zero temperature state, all spins pointing in one direction. Then clusterupdates are performed until over 100k sweeps have been performed.

(We say that one sweep has been performed when a number of spins equal to the total number of spins on the lattice has been tested to be added to a cluster.
Thus the number of clusters and sweeps performed as warmup varies slightly between independent simulations)

After the warmup, cluster updates are performed for 100k more sweeps, and the thermodynamic quantities are collected after each update. Then, averages are calculated and printed. Then 100 sweeps are performed to take the lattice to a state not correlated, and another 100k sample sweeps are performed. This is repeated 100 times. The system sizes simulated are 4, 8, 16, 32, 64 and 128. 
The number of simulations currently used to calculate $U_4, \rho_s, \chi $ etc. for the system sizes 4


\begin{table}[htpb]
\begin{center}
\begin{tabular}{l l l}
  L & No. simulations & No. uncorrelated averages\\
  4 & 110 & 11000\\
  8 & 110 & 11000\\
  16 & 45.54 & 4554\\
  32 & 12 & 1200\\
  64 & 1.16 & 116 \\
  128 & 0  & 0 \\
\end{tabular}
\end{center}
\caption{Number of simulations performed on the 3DXY-model}
\end{table}

\begin{table}[htpb]
\begin{center}
\begin{tabular}{l l l}
  L & No. simulations & No. uncorrelated averages\\
  4 & 100 & 10000\\
  8 & 100 & 10000\\
  16 & 100 & 10000\\
  32 & 97.52 & 9752 \\
  64 & 5.8 & 580 \\
  128 & 1.15  & $115^{*}$ \\
\end{tabular}
\end{center}
\caption{Number of simulations performed on the 3D Ising-model}
\end{table}
\subsection{3D Ising-model}
The same general formula is used when simulating the 3D Ising-model, however, the temperature range is changed to capture the phase transition, $T_{\trm{run}} = 4.50, ~ T_{\trm{range}} = 4.486-4.515$.



%\begin{figure}[!htpb]
  \centering
  \includegraphics[width=15cm]{./plots/3DXY/3DXY_BinderCumulant_vs_Temperature.eps}
  \caption{3DXY}
\end{figure}

\begin{figure}[!htpb]
  \centering
  \includegraphics[width=15cm]{./plots/3DXY/3DXY_Susceptibility_vs_L.eps}
  \caption{3DXY}
\end{figure}

\begin{figure}[!htpb]
  \centering
  \includegraphics[width=15cm]{./plots/3DXY/3DXY_Equilibration_study.eps}
  \caption{3DXY}
\end{figure}

\begin{figure}[!htpb]
  \centering
  \includegraphics[width=15cm]{./plots/3DXY/3DXY_SuperfluidDensity_vs_L.eps}
  \caption{3DXY}
\end{figure}

\begin{figure}[!htpb]
  \centering
  \includegraphics[width=15cm]{./plots/3DXY/3DXY_BinderCumulant_vs_L.eps}
  \caption{3DXY}
\end{figure}

\begin{figure}[!htpb]
  \centering
  \includegraphics[width=15cm]{./plots/3DXY/3DXY_dBdT_vs_L.eps}
  \caption{3DXY}
\end{figure}

\begin{figure}[!htpb]
  \centering
  \includegraphics[width=15cm]{./plots/3DXY/std_intersect_RS_vs_omega.eps}
  \caption{3DXY}
\end{figure}

\begin{figure}[!htpb]
  \centering
  \includegraphics[width=15cm]{./plots/3DXY/3DXY_find_teq_scaling.eps}
  \caption{3DXY}
\end{figure}

\begin{figure}[!htpb]
  \centering
  \includegraphics[width=15cm]{./plots/3DXY/3DXY_SuperfluidDensity_vs_Temperature.eps}
  \caption{3DXY}
\end{figure}

\begin{figure}[!htpb]
  \centering
  \includegraphics[width=15cm]{./plots/3DXY/3DXY_Energy_vs_Temperature.eps}
  \caption{3DXY}
\end{figure}

\begin{figure}[!htpb]
  \centering
  \includegraphics[width=15cm]{./plots/3DXY/3DXY_omega_from_rs_vs_Temperature.eps}
  \caption{3DXY}
\end{figure}

\begin{figure}[!htpb]
  \centering
  \includegraphics[width=15cm]{./plots/3DXY/3DXY_Magnetization_vs_L.eps}
  \caption{3DXY}
\end{figure}

\begin{figure}[!htpb]
  \centering
  \includegraphics[width=15cm]{./plots/3DXY/3DXY_Susceptibility_vs_Temperature.eps}
  \caption{3DXY}
\end{figure}

\begin{figure}[!htpb]
  \centering
  \includegraphics[width=15cm]{./plots/3DXY/3DXY_dBdT_vs_Temperature.eps}
  \caption{3DXY}
\end{figure}

\begin{figure}[!htpb]
  \centering
  \includegraphics[width=15cm]{./plots/3DXY/3DXY_omega_from_B_vs_Temperature.eps}
  \caption{3DXY}
\end{figure}

\begin{figure}[!htpb]
  \centering
  \includegraphics[width=15cm]{./plots/3DXY/3DXY_Magnetization_vs_Temperature.eps}
  \caption{3DXY}
\end{figure}

\begin{figure}[!htpb]
  \centering
  \includegraphics[width=15cm]{./plots/3DXY/3DXY_Energy_vs_L.eps}
  \caption{3DXY}
\end{figure}

\begin{figure}[!htpb]
  \centering
  \includegraphics[width=15cm]{./plots/3DXY/std_intersect_Bin_vs_omega.eps}
  \caption{3DXY}
\end{figure}

\begin{figure}[!htpb]
  \centering
  \includegraphics[width=15cm]{./plots/Ising3D/Ising3D_M2_vs_L.eps}
  \caption{Ising3D}
\end{figure}

\begin{figure}[!htpb]
  \centering
  \includegraphics[width=15cm]{./plots/Ising3D/Ising3D_M4_vs_L.eps}
  \caption{Ising3D}
\end{figure}

\begin{figure}[!htpb]
  \centering
  \includegraphics[width=15cm]{./plots/Ising3D/Ising3D_BinderCumulant_vs_L.eps}
  \caption{Ising3D}
\end{figure}

\begin{figure}[!htpb]
  \centering
  \includegraphics[width=15cm]{./plots/Ising3D/Ising3D_Susceptibility_vs_L.eps}
  \caption{Ising3D}
\end{figure}

\begin{figure}[!htpb]
  \centering
  \includegraphics[width=15cm]{./plots/Ising3D/Ising3D_Susceptibility_vs_Temperature.eps}
  \caption{Ising3D}
\end{figure}

\begin{figure}[!htpb]
  \centering
  \includegraphics[width=15cm]{./plots/Ising3D/Ising3D_M2_vs_Temperature.eps}
  \caption{Ising3D}
\end{figure}

\begin{figure}[!htpb]
  \centering
  \includegraphics[width=15cm]{./plots/Ising3D/Ising3D_Energy_vs_Temperature.eps}
  \caption{Ising3D}
\end{figure}

\begin{figure}[!htpb]
  \centering
  \includegraphics[width=15cm]{./plots/Ising3D/Ising3D_omega_from_B_vs_Temperature.eps}
  \caption{Ising3D}
\end{figure}

\begin{figure}[!htpb]
  \centering
  \includegraphics[width=15cm]{./plots/Ising3D/Ising3D_HeatCapacity_vs_L.eps}
  \caption{Ising3D}
\end{figure}

\begin{figure}[!htpb]
  \centering
  \includegraphics[width=15cm]{./plots/Ising3D/Ising3D_Energy_vs_L.eps}
  \caption{Ising3D}
\end{figure}

\begin{figure}[!htpb]
  \centering
  \includegraphics[width=15cm]{./plots/Ising3D/Ising3D_HeatCapacity_vs_Temperature.eps}
  \caption{Ising3D}
\end{figure}

\begin{figure}[!htpb]
  \centering
  \includegraphics[width=15cm]{./plots/Ising3D/Ising3D_M4_vs_Temperature.eps}
  \caption{Ising3D}
\end{figure}

\begin{figure}[!htpb]
  \centering
  \includegraphics[width=15cm]{./plots/Ising3D/Equilibration time study.eps}
  \caption{Ising3D}
\end{figure}

\begin{figure}[!htpb]
  \centering
  \includegraphics[width=15cm]{./plots/Ising3D/Ising3D_BinderCumulant_vs_Temperature.eps}
  \caption{Ising3D}
\end{figure}

\begin{figure}[!htpb]
  \centering
  \includegraphics[width=15cm]{./plots/Ising3D/std_intersect_Bin_vs_omega.eps}
  \caption{Ising3D}
\end{figure}







\chapter{Summary and conclusions}\label{ch:Summary}
%        File: preface.tex
%     Created: mån feb 12 01:00  2018 C
% Last Change: mån feb 12 01:00  2018 C
%
This is the summary.




\begin{figure}[!htpb]
  \centering
  \includegraphics[width=15cm]{./plots/3DXY/3DXY_BinderCumulant_vs_Temperature.eps}
  \caption{3DXY}
\end{figure}

\begin{figure}[!htpb]
  \centering
  \includegraphics[width=15cm]{./plots/3DXY/3DXY_Susceptibility_vs_L.eps}
  \caption{3DXY}
\end{figure}

\begin{figure}[!htpb]
  \centering
  \includegraphics[width=15cm]{./plots/3DXY/3DXY_Equilibration_study.eps}
  \caption{3DXY}
\end{figure}

\begin{figure}[!htpb]
  \centering
  \includegraphics[width=15cm]{./plots/3DXY/3DXY_SuperfluidDensity_vs_L.eps}
  \caption{3DXY}
\end{figure}

\begin{figure}[!htpb]
  \centering
  \includegraphics[width=15cm]{./plots/3DXY/3DXY_BinderCumulant_vs_L.eps}
  \caption{3DXY}
\end{figure}

\begin{figure}[!htpb]
  \centering
  \includegraphics[width=15cm]{./plots/3DXY/3DXY_dBdT_vs_L.eps}
  \caption{3DXY}
\end{figure}

\begin{figure}[!htpb]
  \centering
  \includegraphics[width=15cm]{./plots/3DXY/std_intersect_RS_vs_omega.eps}
  \caption{3DXY}
\end{figure}

\begin{figure}[!htpb]
  \centering
  \includegraphics[width=15cm]{./plots/3DXY/3DXY_find_teq_scaling.eps}
  \caption{3DXY}
\end{figure}

\begin{figure}[!htpb]
  \centering
  \includegraphics[width=15cm]{./plots/3DXY/3DXY_SuperfluidDensity_vs_Temperature.eps}
  \caption{3DXY}
\end{figure}

\begin{figure}[!htpb]
  \centering
  \includegraphics[width=15cm]{./plots/3DXY/3DXY_Energy_vs_Temperature.eps}
  \caption{3DXY}
\end{figure}

\begin{figure}[!htpb]
  \centering
  \includegraphics[width=15cm]{./plots/3DXY/3DXY_omega_from_rs_vs_Temperature.eps}
  \caption{3DXY}
\end{figure}

\begin{figure}[!htpb]
  \centering
  \includegraphics[width=15cm]{./plots/3DXY/3DXY_Magnetization_vs_L.eps}
  \caption{3DXY}
\end{figure}

\begin{figure}[!htpb]
  \centering
  \includegraphics[width=15cm]{./plots/3DXY/3DXY_Susceptibility_vs_Temperature.eps}
  \caption{3DXY}
\end{figure}

\begin{figure}[!htpb]
  \centering
  \includegraphics[width=15cm]{./plots/3DXY/3DXY_dBdT_vs_Temperature.eps}
  \caption{3DXY}
\end{figure}

\begin{figure}[!htpb]
  \centering
  \includegraphics[width=15cm]{./plots/3DXY/3DXY_omega_from_B_vs_Temperature.eps}
  \caption{3DXY}
\end{figure}

\begin{figure}[!htpb]
  \centering
  \includegraphics[width=15cm]{./plots/3DXY/3DXY_Magnetization_vs_Temperature.eps}
  \caption{3DXY}
\end{figure}

\begin{figure}[!htpb]
  \centering
  \includegraphics[width=15cm]{./plots/3DXY/3DXY_Energy_vs_L.eps}
  \caption{3DXY}
\end{figure}

\begin{figure}[!htpb]
  \centering
  \includegraphics[width=15cm]{./plots/3DXY/std_intersect_Bin_vs_omega.eps}
  \caption{3DXY}
\end{figure}

\begin{figure}[!htpb]
  \centering
  \includegraphics[width=15cm]{./plots/Ising3D/Ising3D_M2_vs_L.eps}
  \caption{Ising3D}
\end{figure}

\begin{figure}[!htpb]
  \centering
  \includegraphics[width=15cm]{./plots/Ising3D/Ising3D_M4_vs_L.eps}
  \caption{Ising3D}
\end{figure}

\begin{figure}[!htpb]
  \centering
  \includegraphics[width=15cm]{./plots/Ising3D/Ising3D_BinderCumulant_vs_L.eps}
  \caption{Ising3D}
\end{figure}

\begin{figure}[!htpb]
  \centering
  \includegraphics[width=15cm]{./plots/Ising3D/Ising3D_Susceptibility_vs_L.eps}
  \caption{Ising3D}
\end{figure}

\begin{figure}[!htpb]
  \centering
  \includegraphics[width=15cm]{./plots/Ising3D/Ising3D_Susceptibility_vs_Temperature.eps}
  \caption{Ising3D}
\end{figure}

\begin{figure}[!htpb]
  \centering
  \includegraphics[width=15cm]{./plots/Ising3D/Ising3D_M2_vs_Temperature.eps}
  \caption{Ising3D}
\end{figure}

\begin{figure}[!htpb]
  \centering
  \includegraphics[width=15cm]{./plots/Ising3D/Ising3D_Energy_vs_Temperature.eps}
  \caption{Ising3D}
\end{figure}

\begin{figure}[!htpb]
  \centering
  \includegraphics[width=15cm]{./plots/Ising3D/Ising3D_omega_from_B_vs_Temperature.eps}
  \caption{Ising3D}
\end{figure}

\begin{figure}[!htpb]
  \centering
  \includegraphics[width=15cm]{./plots/Ising3D/Ising3D_HeatCapacity_vs_L.eps}
  \caption{Ising3D}
\end{figure}

\begin{figure}[!htpb]
  \centering
  \includegraphics[width=15cm]{./plots/Ising3D/Ising3D_Energy_vs_L.eps}
  \caption{Ising3D}
\end{figure}

\begin{figure}[!htpb]
  \centering
  \includegraphics[width=15cm]{./plots/Ising3D/Ising3D_HeatCapacity_vs_Temperature.eps}
  \caption{Ising3D}
\end{figure}

\begin{figure}[!htpb]
  \centering
  \includegraphics[width=15cm]{./plots/Ising3D/Ising3D_M4_vs_Temperature.eps}
  \caption{Ising3D}
\end{figure}

\begin{figure}[!htpb]
  \centering
  \includegraphics[width=15cm]{./plots/Ising3D/Equilibration time study.eps}
  \caption{Ising3D}
\end{figure}

\begin{figure}[!htpb]
  \centering
  \includegraphics[width=15cm]{./plots/Ising3D/Ising3D_BinderCumulant_vs_Temperature.eps}
  \caption{Ising3D}
\end{figure}

\begin{figure}[!htpb]
  \centering
  \includegraphics[width=15cm]{./plots/Ising3D/std_intersect_Bin_vs_omega.eps}
  \caption{Ising3D}
\end{figure}


% This starts the appendices.
%\appendix

%\chapter{Appendix about something}
%\input{appendix}
\addcontentsline{toc}{chapter}{Bibliography}
\bibliographystyle{thesis_bib_style}
\bibliography{references}

%\part{Scientific papers}

%\input{prepapers}

\end{document}





